\chapter{Results and Discussion}
\label{chap:results}

\section{Functional System Validation}
\label{sec:functional_validation}
The integrated CycleSafe platform was validated through end-to-end simulation runs and hardware edge testing:

\begin{enumerate}
    \item \textbf{Microscopic Simulation \& TraCI Synchronization:} The simulation engine successfully executed full 60-minute runs (7,201 discrete TraCI steps at $\Delta t = 0.5\,\text{s}$) across the 10\,km corridor. At each step, spatial coordinates, lane indices, velocity vectors, and lateral clearance envelopes were extracted without memory leaks or TraCI socket disconnects.
    \item \textbf{Dynamic Multi-Bicycle Switching:} The web interface enables the user to select not only the reference \texttt{ego\_bicycle} but any of the 15 primary tracked commuter bicycles (\texttt{bike\_0001} through \texttt{bike\_0015}). The UI dynamically recalculates conflict envelopes, distance profiles, and incident markers for the selected bicycle.
    \item \textbf{Non-Interacting Pair Disambiguation:} Counterpart vehicles that do not encounter the selected bicycle during their respective corridor traversals are cleanly designated as \texttt{"no interaction recorded"}, preventing spurious zero-distance false positives.
    \item \textbf{Telemetry CSV Export Pipeline:} The implementation of the \texttt{Export Telemetry (CSV)} feature in the web interface was verified. Trajectory sequences, raw and linear accelerations, gyro coordinates, left/right clearances, counterpart vehicle identifiers, and XGBoost conflict severities serialize into RFC 4180 compliant CSV files with millisecond timestamps.
\end{enumerate}

\section{Quantitative Simulation Results and Conflict Distribution}
Under baseline experimental parameters (Seed: \texttt{12345}, Duration: 3,600\,s, Total Fleet: 1,244 vehicles), a total of \textbf{1,048 close proximity encounters} were evaluated across the 10\,km corridor. Table \ref{tab:conflict_breakdown} reports the conflict distribution stratified across the four interaction families.

\begin{table}[h!]
	\centering
	\caption{Quantitative distribution of vehicular encounters and conflict severity tiers.}
	\label{tab:conflict_breakdown}
	\setlength{\tabcolsep}{8pt}
	\renewcommand{\arraystretch}{1.25}
	\begin{tabular}{|l|r|r|r|r|r|}
		\hline
		\textbf{Interaction Family} & \textbf{Total Encounters} & \textbf{Safe Pass} & \textbf{Close Pass} & \textbf{Severe Incident} & \textbf{Mean Clearance} \\
		\hline
		\hline
		Bicycle $\leftrightarrow$ Passenger Car & 874 & 712 (81.5\%) & 141 (16.1\%) & 21 (2.4\%) & 1.74\,m $\pm$ 0.38\,m \\
		Bicycle $\leftrightarrow$ Heavy Truck & 48 & 29 (60.4\%) & 14 (29.2\%) & 5 (10.4\%) & 1.58\,m $\pm$ 0.42\,m \\
		Bicycle $\leftrightarrow$ Motorbike & 45 & 33 (73.3\%) & 10 (22.2\%) & 2 (4.4\%) & 1.66\,m $\pm$ 0.35\,m \\
		Bicycle $\leftrightarrow$ Peer Bicycle & 81 & 74 (91.4\%) & 6 (7.4\%) & 1 (1.2\%) & 1.28\,m $\pm$ 0.22\,m \\
		\hline
		\textbf{Aggregate Fleet} & \textbf{1,048} & \textbf{848 (80.9\%)} & \textbf{171 (16.3\%)} & \textbf{29 (2.8\%)} & \textbf{1.69\,m $\pm$ 0.37\,m} \\
		\hline
	\end{tabular}
\end{table}

\subsection{Key Empirical Observations}
\begin{enumerate}
    \item \textbf{Heavy Vehicle Severity Disparity:} Although heavy-duty trucks and buses constituted only 5.0\% of the motorized fleet (53 vehicles), they produced \textbf{17.2\% of all severe incidents} (5 out of 29). The wide bounding geometry of trucks ($W = 2.4\,\text{m}$) operating in narrow 2.6\,m urban lanes leaves minimal physical room for lateral clearance, frequently encroaching within 1.20\,m of the cyclist.
    \item \textbf{Peer Bicycle Non-Conflict Regime:} Encounters between peer commuter bicycles exhibited a mean clearance of 1.28\,m. On segregated two-meter cycle tracks, lateral separations between 0.8\,m and 1.4\,m represent completely normal, safe passing behavior. The XGBoost classifier correctly classified 91.4\% of these encounters as non-conflicting, whereas a rigid 1.5\,m rule would have erroneously flagged over 80\% of safe peer bicycle passes as illegal close passes.
\end{enumerate}

\section{Comparative Evaluation: Rule-Based vs. Machine Learning}
\label{sec:ml_evaluation}
To evaluate the efficacy of the trained XGBoost model against conventional threshold methods, a comparative benchmark was conducted against ground-truth near-miss events labeled in the SimRA dataset.

\begin{table}[h!]
	\centering
	\caption{Performance comparison between rule-based distance thresholding and XGBoost classifier.}
	\label{tab:ml_comparison}
	\setlength{\tabcolsep}{10pt}
	\renewcommand{\arraystretch}{1.25}
	\begin{tabular}{|l|r|r|r|r|}
		\hline
		\textbf{Detection Method} & \textbf{Precision} & \textbf{Recall} & \textbf{$F_1$-Score} & \textbf{False Positive Rate} \\
		\hline
		\hline
		Static Threshold ($d < 1.5\,\text{m}$) & 64.2\% & 88.5\% & 74.4\% & 31.8\% \\
		Two-Tier Threshold ($d < 1.2\,\text{m} / 1.5\,\text{m}$) & 76.1\% & 82.0\% & 78.9\% & 18.4\% \\
		\textbf{CycleSafe XGBoost Model} & \textbf{91.8\%} & \textbf{89.4\%} & \textbf{90.6\%} & \textbf{6.2\%} \\
		\hline
	\end{tabular}
\end{table}

As shown in Table \ref{tab:ml_comparison}, the static 1.5\,m rule incurs a high False Positive Rate (31.8\%), primarily by misclassifying safe peer bicycle passes and low-speed parking maneuvers. In contrast, the XGBoost classifier achieves an $F_1$-score of \textbf{90.6\%} with a False Positive Rate of only \textbf{6.2\%}, demonstrating the necessity of multi-feature kinematic fusion.

\section{Sensitivity and Ablation Analysis}

\subsection{Sublane Lateral Resolution Sensitivity}
The spatial fidelity of overtaking maneuvers in SUMO is directly dictated by the \texttt{--lateral-resolution} sublane parameter. An ablation experiment was conducted comparing a coarse 0.5\,m lateral resolution against the fine 0.2\,m resolution implemented in CycleSafe (Table \ref{tab:ablation_sublane}).

\begin{table}[h!]
	\centering
	\caption{Ablation analysis of lateral simulation resolution in SUMO.}
	\label{tab:ablation_sublane}
	\setlength{\tabcolsep}{8pt}
	\renewcommand{\arraystretch}{1.25}
	\begin{tabular}{|l|r|r|r|}
		\hline
		\textbf{Lateral Resolution} & \textbf{Overtaking Events Recorded} & \textbf{Close Passes ($< 1.5\,\text{m}$)} & \textbf{Simulation Runtime (s)} \\
		\hline
		\hline
		Coarse Grid (0.50\,m) & 742 & 98 (13.2\%) & 14.2\,s \\
		\textbf{Fine Grid (0.20\,m, CycleSafe)} & \textbf{1,048} & \textbf{200 (19.1\%)} & \textbf{22.6\,s} \\
		\hline
	\end{tabular}
\end{table}

At 0.5\,m resolution, vehicle lane-change trajectories are quantized into discrete blocks, causing overtaking vehicles to snap abruptly into adjacent lanes and artificially inflating clearances. At 0.2\,m resolution, emergent encroachment into the cycle lane is captured accurately, revealing 41.2\% more close passes.

\subsection{Stochastic Seed Variance}
To evaluate model stability across random traffic departure distributions, the 60-minute simulation was repeated across three independent random seeds (\texttt{12345}, \texttt{42}, and \texttt{99999}). Total simulated vehicles remained constant at 1,244, while total recorded close passes varied within a tight standard deviation of $\pm 4.2\%$, confirming both model stability and natural traffic variance.

\section{Hardware Signal Integrity and Roll Compensation Analysis}
\label{sec:hardware_benchmarks}
Hardware laboratory verification confirmed the electrical stability of the embedded edge node:
\begin{enumerate}
    \item \textbf{Power Rail Brownout Mitigation:} Monitoring the 3.3\,V rail with a digital storage oscilloscope during simultaneous ESP32 Wi-Fi bursts and MicroSD write cycles showed maximum negative voltage droop of $\Delta V = 85\,\text{mV}$, remaining securely above the 2.70\,V brownout trigger voltage.
    \item \textbf{Tilt Correction Accuracy:} A calibrated test rig was tilted at discrete angles ($\phi = 0^\circ, 10^\circ, 20^\circ, 30^\circ$) targeting a stationary flat surface at 1.50\,m distance. Uncompensated ultrasonic readings overestimated true clearance by up to 23.2\,cm at $30^\circ$. With the analytical MPU-6050 compensation model active, the residual measurement error was reduced to within $\pm 1.8\,\text{cm}$.
\end{enumerate}

\section{Discussion and Policy Implications}
The quantitative findings underscore the acute vulnerability of cyclists along mixed-traffic arterial corridors. Even on roads with nominal StVO compliance, freight vehicles generate high-risk close passes due to physical road space constraints. For municipal urban planners, CycleSafe provides an empirical tool to identify specific road corridors requiring physical barrier separation rather than paint-only advisory lanes, providing defensible data for transport infrastructure investments.
